\section*{Capítulo \ref{ch:b2:planet-formation-meteorites} --- Formação dos planetas e meteoritos}

\begin{solution}{exo:b2:planet-formation-meteorites:1}
Gás: \qty{2200}{g.cm^{-2}} a \qty{1}{au}; a \qty{5.2}{au},
$2200\times5.2^{-3/2} = 2200/11.86 = \qty{186}{g.cm^{-2}}$. Sólidos:
$33\times0.33 = \qty{10.9}{g.cm^{-2}}$ a \qty{1}{au};
$33/11.86 = \qty{2.78}{g.cm^{-2}}$ a \qty{5.2}{au}. A \qty{5.2}{au}, o gás é
$186/2.78 = 67$ vezes mais pesado que os sólidos.
\end{solution}

\begin{solution}{exo:b2:planet-formation-meteorites:2}
Alumínio (\qty{1653}{K}, ultrarrefratário), magnésio (\qty{1336}{K}, comum),
ferro (\qty{1334}{K}, comum), sódio (\qty{958}{K}, moderadamente volátil), enxofre
(\qty{664}{K}, volátil), oxigênio como gelo de água (\qty{180}{K}, altamente volátil).
As CAIs, os primeiros condensados, são enriquecidas em alumínio (e em cálcio e titânio);
a Terra é empobrecida nos voláteis, sódio, enxofre e água, em relação ao
magnésio e ao alumínio.
\end{solution}

\begin{solution}{exo:b2:planet-formation-meteorites:3}
\omterm{def:b2:planet-formation-meteorites:fall-find}{Quedas}: \omterm{def:b2:planet-formation-meteorites:chondrule}{condritos} ordinários $856/1107 = \qty{77.3}{\percent}$, ferrosos
$49/1107 = \qty{4.4}{\percent}$. \omterm{def:b2:planet-formation-meteorites:fall-find}{Achados}: $38647/44534 = \qty{86.8}{\percent}$ e
$1020/44534 = \qty{2.3}{\percent}$. \omterm{def:b2:planet-formation-meteorites:fall-find}{Achados} antárticos: $22096/44534 =
\qty{49.6}{\percent}$. As \omterm{def:b2:planet-formation-meteorites:fall-find}{quedas}: toda \omterm{def:b2:planet-formation-meteorites:fall-find}{queda} observada é coletada, qualquer que seja sua
classe, ao passo que um \omterm{def:b2:planet-formation-meteorites:fall-find}{achado} depende de quão bem o \omterm{def:b2:planet-formation-meteorites:fall-find}{meteorito} sobrevive no solo,
de sua aparência e de onde as pessoas procuram, e metade dos \omterm{def:b2:planet-formation-meteorites:fall-find}{achados} vem de um único tipo de
lugar.
\end{solution}

\begin{solution}{exo:b2:planet-formation-meteorites:4}
$\Delta t = 1.035\ln(\num{5.23e-5}/\num{1.0e-5}) = 1.035\ln 5.23 =
1.035\times1.654 = \qty{1.71}{Myr}$ depois das CAIs.
\end{solution}

\begin{solution}{exo:b2:planet-formation-meteorites:5}
$4\pi\Sigma_1 r_1^2 = 4\pi\times\num{2.2e4}\times(\num{1.496e11})^2 =
\qty{6.19e27}{kg}$. Dentro de \qty{30}{au}: $\times\sqrt{30} = 5.48$, portanto
\qty{3.39e28}{kg} $= 0.017$ massa solar; dentro de \qty{100}{au}: $\times10$,
\qty{6.19e28}{kg} $= 0.031$ massa solar. Ambas ficam dentro da faixa observada,
mais perto de seu limite superior: a nebulosa era um disco comum.
\end{solution}

\begin{solution}{exo:b2:planet-formation-meteorites:6}
$r = (120/T)^{7/3}$ au: para a água, $(120/182)^{7/3} = 0.659^{7/3} =
\qty{0.38}{au}$; para o metano, $(120/41)^{7/3} = 2.93^{7/3} = \qty{12.3}{au}$.
O grão iluminado diretamente pelo Sol é a camada superficial do disco; o plano médio,
protegido pela poeira acima dele, é mais frio, de modo que o gelo sobrevive mais perto. Os
\omterm{def:b2:planet-formation-meteorites:planetesimal}{planetesimais} se formaram no plano médio, portanto vale a estimativa do plano médio, desde que
o disco seja passivo; o aquecimento por acreção desloca a \omterm{def:b2:planet-formation-meteorites:snow-line}{linha de neve} para fora.
\end{solution}

\begin{solution}{exo:b2:planet-formation-meteorites:7}
A \qty{1}{au}: $r^2 = \num{2.238e22}$, $\Sigma_p = \qty{109}{kg.m^{-2}}$, portanto
$2\pi b r^2\Sigma_p = 31.42\times\num{2.238e22}\times109 = \num{7.66e25}$;
$(\num{7.66e25})^{3/2} = \num{6.71e38}$, dividido por
$(3\times\num{1.989e30})^{1/2} = \num{2.443e15}$: \qty{2.75e23}{kg}
$= 0.046$ massa terrestre. A \qty{5.2}{au}: $r^2 = \num{6.05e23}$,
$\Sigma_p = \qty{27.8}{kg.m^{-2}}$, $2\pi b r^2\Sigma_p = \num{5.29e26}$,
$M = \num{1.22e40}/\num{2.443e15} = \qty{4.98e24}{kg} = 0.83$ massa terrestre.
Com \num{3.5} vezes mais sólidos, $M$ cresce como $3.5^{3/2} = 6.55$: 5.5 massas
terrestres.
\end{solution}

\begin{solution}{exo:b2:planet-formation-meteorites:8}
Por quilograma, $\num{1.01e19}$ átomos, portanto
$\num{1.01e19}\times\num{5.00e-13} = \qty{5.05e6}{J.kg^{-1}}$. O aumento de
temperatura seria $\num{5.05e6}/837 = \qty{6030}{K}$. Aquecer de \num{250} a
\qty{1700}{K} exige $837\times1450 = \qty{1.21e6}{J.kg^{-1}}$: o alumínio-26
fornecia cerca de quatro vezes isso.
\end{solution}

\begin{solution}{exo:b2:planet-formation-meteorites:9}
$\num{5.07e10}/\num{7.49e7} = 677$. A vida média do ferro-60 é
$2.62/\ln2 = \qty{3.78}{Myr}$. Potências iniciais por grama:
$\num{5.07e10}/1.035 = \num{4.90e10}$ e $\num{7.49e7}/3.78 = \num{1.98e7}$ erg por
grama por milhão de anos, uma razão de cerca de \num{2470}: o alumínio-26 dominava de longe.
\end{solution}

\begin{solution}{exo:b2:planet-formation-meteorites:10}
As potências são iguais quando $\num{4.90e10}\,e^{-t/1.035} = \num{1.98e7}\,e^{-t/3.78}$,
isto é, $t(1/1.035 - 1/3.78) = \ln(\num{4.90e10}/\num{1.98e7}) = \ln 2475 = 7.81$;
com $0.966 - 0.265 = 0.702$ por milhão de anos, $t = \qty{11.1}{Myr}$. Nesse momento, o
ferro-60 tem ainda $\num{7.49e7}\,e^{-11.1/3.78} = \num{7.49e7}\times0.053 =
\qty{4.0e6}{erg.g^{-1}}$ a liberar, isto é, \qty{400}{J.kg^{-1}}: milhares de vezes
menos do que o necessário para fundir qualquer coisa.
\end{solution}

\begin{solution}{exo:b2:planet-formation-meteorites:11}
$\tau = 1.035\times\num{3.156e13} = \qty{3.267e13}{s}$ e
$H_0\tau = \qty{6.21e6}{J.kg^{-1}}$. Para \qty{1200}{K}:
$t_a^* = 1.035\ln\bigl(\num{6.21e6}/(837\times950)\bigr) = 1.035\ln7.81 =
\qty{2.13}{Myr}$; para \qty{1700}{K}: $1.035\ln\bigl(\num{6.21e6}/(837\times1450)\bigr)
= 1.035\ln5.12 = \qty{1.69}{Myr}$. Reduzir $H_0$ à metade subtrai $\tau\ln2 =
\qty{0.72}{Myr}$ dos dois, pois $t_a^* = \tau\ln H_0 + \tau\ln(\tau/c(T_c -
T_0))$ e só o primeiro termo muda: o deslocamento é uma meia-vida, qualquer que seja a
temperatura crítica.
\end{solution}

\begin{solution}{exo:b2:planet-formation-meteorites:12}
$\kappa = 2.1/(3300\times837) = \qty{7.60e-7}{m^2.s^{-1}}$ e
$\sqrt{\kappa\tau} = \sqrt{\num{7.60e-7}\times\num{3.267e13}} = \qty{4980}{m}$,
cerca de \qty{5}{km}. O centro de um corpo de \qty{100}{km} de diâmetro está a \qty{50}{km} de profundidade,
dez vezes essa distância: seu calor não consegue escapar enquanto o alumínio-26 decai, e a
proposição se aplica. Um corpo de \qty{2}{km} de diâmetro perde seu calor tão depressa quanto o
produz. A fusão exige um raio de algumas vezes \qty{5}{km}, isto é, um diâmetro de uns
\qtyrange{10}{20}{km} ou mais, como encontram os modelos.
\end{solution}

\begin{solution}{pb:b2:planet-formation-meteorites:1}
\textbf{1.} Uma \omterm{def:b2:planet-formation-meteorites:fall-find}{queda}: sua chegada foi vista e ela foi recolhida na manhã seguinte.
Ela não está alterada pelo intemperismo e se soma à amostra não enviesada do que chega à Terra.

\textbf{2.} $856/1107 = 0.77$: cerca de três chances em quatro.

\textbf{3.} \omterm{def:b2:planet-formation-meteorites:chondrule}{Côndrulos} de \qty{0.3}{mm} de diâmetro, \qty{9}{\percent} de metal e olivina com
\qty{18}{\percent} de faialita: um \omterm{def:b2:planet-formation-meteorites:chondrule}{condrito} H (\qty{10}{\percent} de metal, Fa 16--20).

\textbf{4.} \omterm{def:b2:planet-formation-meteorites:chondrule}{Côndrulos} esmaecidos e minerais uniformes indicam metamorfismo térmico: tipos
5 a 6. Os \omterm{def:b2:planet-formation-meteorites:chondrule}{condritos} ordinários atingiram no máximo cerca de \qty{1220}{K}, isto é,
\qty{947}{\degreeCelsius}.

\textbf{5.} Os \omterm{def:b2:planet-formation-meteorites:chondrule}{condritos} ordinários contêm apenas \qtyrange{0.1}{1}{\percent} de CAIs em
volume. A pedra é feita de \omterm{def:b2:planet-formation-meteorites:chondrule}{côndrulos} e matriz agregados depois das CAIs, que
são os sólidos mais antigos: toda rocha que os contém é mais jovem.

\textbf{6.} Cada mineral recebeu na formação o mesmo magnésio mais alumínio-26
em proporção ao seu alumínio; desde então, o alumínio-26 transformou-se em
magnésio-26, de modo que, quanto mais alumínio por magnésio um mineral tinha, mais
magnésio-26 ele ganhou.

\textbf{7.} $(0.139590 - 0.139500)/(15.0 - 0.05) = \num{0.000090}/14.95 =
\num{6.02e-6}$: a razão $\ce{^{26}Al}/\ce{^{27}Al}$ do \omterm{def:b2:planet-formation-meteorites:chondrule}{côndrulo} quando ele
se formou.

\textbf{8.} $0.139500 + \num{6.02e-6}\times3.00 = 0.139518$: sobre a reta.

\textbf{9.} $\Delta t = 1.035\ln(\num{5.23e-5}/\num{6.02e-6}) = 1.035\ln8.69 =
1.035\times2.162 = \qty{2.24}{Myr}$.

\textbf{10.} $4567.5 - 2.2 = \qty{4565.3}{\Ma}$.

\textbf{11.} $\num{5.23e-5}\,e^{-10/1.035} = \num{5.23e-5}\times\num{6.4e-5} =
\num{3.3e-9}$, magnésio-26 pouco demais para ser medido. O háfnio-182 (meia-vida de
\qty{8.9}{Myr}) ou o manganês-53 (\qty{3.7}{Myr}) convêm a acontecimentos ocorridos uns dez milhões de anos
depois das CAIs.

\textbf{12.} $\kappa = 2.1/(3300\times837) = \qty{7.60e-7}{m^2.s^{-1}}$;
$R^2/\kappa = \num{4e8}/\num{7.60e-7} = \qty{5.26e14}{s} = \qty{16.7}{Myr}$, dezesseis
vezes a vida média do alumínio-26: o centro conserva seu calor.

\textbf{13.} $\num{1.9e-7}\,e^{-2.4/1.035} = \num{1.9e-7}\times0.0984 =
\qty{1.87e-8}{W.kg^{-1}}$.

\textbf{14.} $\num{1.87e-8}\times\num{3.267e13} = \qty{6.11e5}{J.kg^{-1}}$.

\textbf{15.} $\Delta T = \num{6.11e5}/837 = \qty{730}{K}$; máximo de
$250 + 730 = \qty{980}{K}$.

\textbf{16.} Acima de \qty{750}{K}, portanto a rocha sinterizou e foi metamorfizada; abaixo de
\qty{1200}{K}, portanto nem o metal nem o sulfeto fundiram. A pedra é um \omterm{def:b2:planet-formation-meteorites:chondrule}{condrito} aquecido,
de \omterm{def:b2:planet-formation-meteorites:type}{tipo petrológico} alto, abaixo dos \qty{1220}{K} da pergunta 4: coerente.

\textbf{17.} $1.035\ln\bigl(\num{6.21e6}/(837\times950)\bigr) = \qty{2.13}{Myr}$:
cerca de \qty{0.3}{Myr} antes.

\textbf{18.} O aumento é $\Delta T(1 - e^{-(t - t_a)/\tau})$; \qty{90}{\percent}
exigem $e^{-(t - t_a)/\tau} = 0.1$, $t - t_a = 1.035\ln10 = \qty{2.38}{Myr}$, isto é,
\qty{4.8}{Myr} depois das CAIs.

\textbf{19.} $\Delta T = 7420\,e^{-0.8/1.035} = 7420\times0.462 = \qty{3430}{K}$,
máximo de \qty{3680}{K}, muito acima dos \qty{1700}{K} em que os silicatos estão fundidos.
O corpo fundiu (o calor de fusão manteve sua temperatura mais baixa) e
diferenciou-se em um núcleo metálico e um manto e uma crosta silicáticos.

\textbf{20.} $e^{(2.4 - 0.8)/1.035} = e^{1.546} = 4.7$ vezes mais.

\textbf{21.} O zinco e o chumbo se condensam perto de \qty{727}{K}, moderadamente voláteis; o magnésio
(\qty{1336}{K}) e o alumínio (\qty{1653}{K}) se condensam muito antes. Os sólidos da
região da pedra se reuniram antes que o gás tivesse esfriado o bastante para que todo o zinco e
todo o chumbo se condensassem.

\textbf{22.} $33\times0.33\times2.5^{-3/2} = 10.89/3.95 = \qty{2.76}{g.cm^{-2}}$,
\qty{27.6}{kg.m^{-2}}.

\textbf{23.} $r^2 = (3.74\times10^{11})^2 = \num{1.40e23}$;
$2\pi\times5\times\num{1.40e23}\times27.6 = \num{1.21e26}$;
$(\num{1.21e26})^{3/2}/\num{2.443e15} = \qty{5.45e23}{kg}$, $0.091$ massa terrestre.
O corpo parental, $\tfrac43\pi(\num{2e4})^3\times3300 = \qty{1.1e17}{kg}$, é cinco
milhões de vezes mais leve: um \omterm{def:b2:planet-formation-meteorites:planetesimal}{planetesimal}, não um embrião.

\textbf{24.} O centro do corpo parental atingiu um máximo de cerca de
\textbf{\qty[text-series-to-math=true]{707}{\degreeCelsius}} ($980 - 273$): quente o bastante para
recristalizar o \omterm{def:b2:planet-formation-meteorites:chondrule}{condrito}, frio demais para fundir seu metal.
\end{solution}
